% GENERATED FILE. Edit canonical lessons, then run npm run generate:books.
% canonical-source-hash: fnv1a64:278dfa6470098905
% canonical-lessons: RU-C489-obsledovat, RU-C489-obsledovaniye, RU-C489-analiz, RU-C489-rentgen, RU-C489-khirurg

\chapter{A Check-up}
\label{ch:ru-a-check-up}
\hlchaptermodality{489}

\begin{chapteropening}
\textbf{By the end of this chapter:} \emph{I can talk about a medical check-up and tests.}

Complete the last lesson of chapter 489: I can talk about a medical check-up and tests.
\end{chapteropening}

\section[\texorpdfstring{\ru{обследовать} (obslédovat)}{обследовать (obslédovat)}]{\ru{обследовать} (obslédovat) — to examine (medically)}

\label{lesson:RU-C489-obsledovat}



\begin{quote}
Before the new one: say the Russian for a nightmare, then the Russian for to doze.
\end{quote}

\subsection*{You'll want to know: \ru{обследовать}}
\textbf{\ru{обследовать}} (\emph{obslédovat}) — \textquotedblleft{}to examine (medically)\textquotedblright{}.

\begin{quote}
\ru{Врач} \ru{внимательно} \ru{обследовал} \ru{пациента}.
\end{quote}

\begin{quote}
\emph{Vrach vnimátelno obslédoval patsiyénta.}
\end{quote}

\begin{quote}
The doctor examined the patient carefully.
\end{quote}

Imperfective and perfective (one form for both). Present \ru{обследую}.

\begin{cousinweb}
\emph{\ru{обследовать}} is built on \emph{\ru{об}-} and \emph{\ru{следовать}} 'to follow', from \emph{\ru{след}} 'track, trail' (EtymDB lists all three without a source).
\end{cousinweb}

\begin{grammarlens}[title={What you've built}]
One more word for telling, explaining and describing at B1.
\end{grammarlens}

\subsection*{Your turn}
\begin{itemize}
\raggedright
  \item \emph{Say it:} \emph{obslédovat}
  \item \emph{Say it:} \emph{obslédovat}, once more
  \item \emph{Say it:} a sentence of your own with \emph{obslédovat}
  \item \emph{Recall:} say the Russian for a nightmare, then the Russian for to doze, then say \emph{obslédovat} again
\end{itemize}

\begin{tcolorbox}[breakable,colback=teal!4,colframe=teal!35!black,title={Before you move on}]
What does \ru{обследовать} mean? (\textquotedblleft{}To examine (medically)\textquotedblright{}.) Say it once more.
\end{tcolorbox}
\section[\texorpdfstring{\ru{обследование} (obslédovaniye)}{обследование (obslédovaniye)}]{\ru{обследование} (obslédovaniye) — a medical examination, a check-up}

\label{lesson:RU-C489-obsledovaniye}



\begin{quote}
Before the new one: say the Russian for to doze, then the Russian for to examine.
\end{quote}

\subsection*{You'll want to know: \ru{обследование}}
\textbf{\ru{обследование}} (\emph{obslédovaniye}) — \textquotedblleft{}a medical examination, a check-up\textquotedblright{}.

\begin{quote}
\ru{Каждый} \ru{год} \ru{я} \ru{прохожу} \ru{обследование}.
\end{quote}

\begin{quote}
\emph{Kázhdyy god ya prokhozhú obslédovaniye.}
\end{quote}

\begin{quote}
Every year I have a check-up.
\end{quote}

Neuter. \emph{\ru{Проходить} \ru{обследование}} — to have a check-up.

\begin{cousinweb}
\emph{\ru{обследование}} is built on \emph{\ru{обследовать}} 'to examine' with \emph{-\ru{ание}}; the root is that of \emph{\ru{след}} 'track, trail', which EtymDB lists without a source, and EtymDB has no entry for the noun.
\end{cousinweb}

\begin{grammarlens}[title={What you've built}]
One more word for telling, explaining and describing at B1.
\end{grammarlens}

\subsection*{Your turn}
\begin{itemize}
\raggedright
  \item \emph{Say it:} \emph{obslédovaniye}
  \item \emph{Say it:} \emph{obslédovaniye}, once more
  \item \emph{Say it:} a sentence of your own with \emph{obslédovaniye}
  \item \emph{Recall:} say the Russian for to doze, then the Russian for to examine, then say \emph{obslédovaniye} again
\end{itemize}

\begin{tcolorbox}[breakable,colback=teal!4,colframe=teal!35!black,title={Before you move on}]
What does \ru{обследование} mean? (\textquotedblleft{}A medical examination, a check-up\textquotedblright{}.) Say it once more.
\end{tcolorbox}
\section[\texorpdfstring{\ru{анализ} (análiz)}{анализ (análiz)}]{\ru{анализ} (análiz) — a medical test; analysis}

\label{lesson:RU-C489-analiz}



\begin{quote}
Before the new one: say the Russian for to examine, then the Russian for a medical examination.
\end{quote}

\subsection*{You'll want to know: \ru{анализ}}
\textbf{\ru{анализ}} (\emph{análiz}) — \textquotedblleft{}a medical test; analysis\textquotedblright{}.

\begin{quote}
\ru{Медсестра} \ru{взяла} \ru{у} \ru{меня} \ru{анализ} \ru{крови}.
\end{quote}

\begin{quote}
\emph{Medsestrá vzyalá u menyá análiz króvi.}
\end{quote}

\begin{quote}
The nurse took a blood sample from me.
\end{quote}

Masculine. \emph{\ru{Сдать} \ru{анализы}} — to have tests done.

\begin{cousinweb}
\emph{\ru{анализ}} is a learned borrowing, usually traced through French or German to Greek \emph{análysis} 'a loosening'; EtymDB has no entry for the Russian word.
\end{cousinweb}

\begin{grammarlens}[title={What you've built}]
One more word for telling, explaining and describing at B1.
\end{grammarlens}

\subsection*{Your turn}
\begin{itemize}
\raggedright
  \item \emph{Say it:} \emph{análiz}
  \item \emph{Say it:} \emph{análiz}, once more
  \item \emph{Say it:} a sentence of your own with \emph{análiz}
  \item \emph{Recall:} say the Russian for to examine, then the Russian for a medical examination, then say \emph{análiz} again
\end{itemize}

\begin{tcolorbox}[breakable,colback=teal!4,colframe=teal!35!black,title={Before you move on}]
What does \ru{анализ} mean? (\textquotedblleft{}A medical test; analysis\textquotedblright{}.) Say it once more.
\end{tcolorbox}
\section[\texorpdfstring{\ru{рентген} (rentgén)}{рентген (rentgén)}]{\ru{рентген} (rentgén) — an X-ray}

\label{lesson:RU-C489-rentgen}



\begin{quote}
Before the new one: say the Russian for a medical examination, then the Russian for a medical test.
\end{quote}

\subsection*{You'll want to know: \ru{рентген}}
\textbf{\ru{рентген}} (\emph{rentgén}) — \textquotedblleft{}an X-ray\textquotedblright{}.

\begin{quote}
\ru{Мне} \ru{сделали} \ru{рентген} \ru{руки}.
\end{quote}

\begin{quote}
\emph{Mne sdélali rentgén rukí.}
\end{quote}

\begin{quote}
They X-rayed my arm.
\end{quote}

Masculine. \emph{\ru{Сделать} \ru{рентген}} — to have an X-ray.

\begin{cousinweb}
\emph{\ru{рентген}} was borrowed from German \emph{Röntgen} 'X-ray' (EtymDB), a word made from the surname of the physicist who found the rays.
\end{cousinweb}

\begin{grammarlens}[title={What you've built}]
One more word for telling, explaining and describing at B1.
\end{grammarlens}

\subsection*{Your turn}
\begin{itemize}
\raggedright
  \item \emph{Say it:} \emph{rentgén}
  \item \emph{Say it:} \emph{rentgén}, once more
  \item \emph{Say it:} a sentence of your own with \emph{rentgén}
  \item \emph{Recall:} say the Russian for a medical examination, then the Russian for a medical test, then say \emph{rentgén} again
\end{itemize}

\begin{tcolorbox}[breakable,colback=teal!4,colframe=teal!35!black,title={Before you move on}]
What does \ru{рентген} mean? (\textquotedblleft{}An X-ray\textquotedblright{}.) Say it once more.
\end{tcolorbox}
\section[\texorpdfstring{\ru{хирург} (khirúrg)}{хирург (khirúrg)}]{\ru{хирург} (khirúrg) — a surgeon}

\label{lesson:RU-C489-khirurg}



\begin{quote}
Before the new one: say the Russian for a medical test, then the Russian for an X-ray.
\end{quote}

\subsection*{You'll want to know: \ru{хирург}}
\textbf{\ru{хирург}} (\emph{khirúrg}) — \textquotedblleft{}a surgeon\textquotedblright{}.

\begin{quote}
\ru{Хирург} \ru{сказал}, \ru{что} \ru{операция} \ru{прошла} \ru{хорошо}.
\end{quote}

\begin{quote}
\emph{Khirúrg skazál, chto operátsiya proshlá khoroshó.}
\end{quote}

\begin{quote}
The surgeon said that the operation had gone well.
\end{quote}

Masculine; also of a woman.

\begin{cousinweb}
\emph{\ru{хирург}} was borrowed from Greek \emph{kheirourgós} 'surgeon' (EtymDB), a compound of \emph{kheír} 'hand' and \emph{érgon} 'work'.
\end{cousinweb}

\begin{grammarlens}[title={What you've built}]
One more word for telling, explaining and describing at B1.
\end{grammarlens}

\subsection*{Your turn}
\begin{itemize}
\raggedright
  \item \emph{Say it:} \emph{khirúrg}
  \item \emph{Say it:} \emph{khirúrg}, once more
  \item \emph{Say it:} a sentence of your own with \emph{khirúrg}
  \item \emph{Recall:} say the Russian for a medical test, then the Russian for an X-ray, then say \emph{khirúrg} again
\end{itemize}

\begin{tcolorbox}[breakable,colback=teal!4,colframe=teal!35!black,title={Before you move on}]
What does \ru{хирург} mean? (\textquotedblleft{}A surgeon\textquotedblright{}.) Say it once more.
\end{tcolorbox}
